% Related Work for Malleable World Models (MWM)
% Typical placement: after % Introduction for Malleable World Models (MWM)
% % Introduction for Malleable World Models (MWM)
% % Introduction for Malleable World Models (MWM)
% \input{introduction}

\section{Introduction}
\label{sec:introduction}

World models enable model-based agents to plan and act by predicting latent dynamics from experience~\cite{hafner2023mastering,hafner2020dream}.
Once trained, however, they are typically deployed with \emph{frozen} parameters.
In open-ended settings---wear on actuators, payload changes, or simulator-to-real gaps---the transition kernel internalized during pretraining ceases to match the environment the agent inhabits.
Prediction error rises, imagined rollouts become misleading, and control performance degrades.
\emph{Test-time adaptation} (TTA) seeks to recover competence after such shifts without returning to expensive offline training.

The dominant TTA recipe is gradient-based learning at deployment: collect post-shift trajectories, backpropagate a surrogate loss, and update network weights~\cite{nagabandi2018learning,zintgraf2021fast}.
This \emph{slow} pathway is expressive but incurs costs that matter in interactive control---batched updates, autograd through long horizons, and tuning to avoid catastrophic forgetting.
When the shift is sudden, the agent may act suboptimally for many steps before the first weight update.

Biological systems often separate a fast, local error-driven store from slower structural learning~\cite{mcclelland1995complementary}.
The fast store is not globally optimal; it provides coarse correction so behavior is not blind during transients.
We transpose this idea to world models: adaptation is not only a motor reflex, but a change in \emph{what the agent believes will happen next}.

We propose \textbf{Malleable World Models (MWM)}: frozen pretrained agents augmented with a \textbf{fast plastic prior residual} $P_t$, updated every step by a Hebbian rule without backprop through pretraining weights, plus optional \textbf{slow} actor and world-model updates from replay.
We evaluate on proprioceptive \texttt{walker-walk} under actuator gain shift ($1.0 \rightarrow 0.7 \rightarrow 1.0$).
Section~\ref{sec:method} details the method; Section~\ref{sec:experiments} reports detection, adaptation latency, and consolidation behavior.

\paragraph{Contributions.}
\begin{enumerate}[leftmargin=*, itemsep=2pt]
  \item A backpropagation-free plastic prior on RSSM logits with norm-capped state $P_t$.
  \item A tri-timescale test-time stack: fast $P_t$, slow actor fine-tuning, and optional WM consolidation with coupled flush.
  \item Empirical characterization on gain shift: rapid trace growth and prior-error changes; control gains when fast plastic scaffolds slow actor learning; pilot evidence for $P_t$ drainage under consolidation.
\end{enumerate}


\section{Introduction}
\label{sec:introduction}

World models enable model-based agents to plan and act by predicting latent dynamics from experience~\cite{hafner2023mastering,hafner2020dream}.
Once trained, however, they are typically deployed with \emph{frozen} parameters.
In open-ended settings---wear on actuators, payload changes, or simulator-to-real gaps---the transition kernel internalized during pretraining ceases to match the environment the agent inhabits.
Prediction error rises, imagined rollouts become misleading, and control performance degrades.
\emph{Test-time adaptation} (TTA) seeks to recover competence after such shifts without returning to expensive offline training.

The dominant TTA recipe is gradient-based learning at deployment: collect post-shift trajectories, backpropagate a surrogate loss, and update network weights~\cite{nagabandi2018learning,zintgraf2021fast}.
This \emph{slow} pathway is expressive but incurs costs that matter in interactive control---batched updates, autograd through long horizons, and tuning to avoid catastrophic forgetting.
When the shift is sudden, the agent may act suboptimally for many steps before the first weight update.

Biological systems often separate a fast, local error-driven store from slower structural learning~\cite{mcclelland1995complementary}.
The fast store is not globally optimal; it provides coarse correction so behavior is not blind during transients.
We transpose this idea to world models: adaptation is not only a motor reflex, but a change in \emph{what the agent believes will happen next}.

We propose \textbf{Malleable World Models (MWM)}: frozen pretrained agents augmented with a \textbf{fast plastic prior residual} $P_t$, updated every step by a Hebbian rule without backprop through pretraining weights, plus optional \textbf{slow} actor and world-model updates from replay.
We evaluate on proprioceptive \texttt{walker-walk} under actuator gain shift ($1.0 \rightarrow 0.7 \rightarrow 1.0$).
Section~\ref{sec:method} details the method; Section~\ref{sec:experiments} reports detection, adaptation latency, and consolidation behavior.

\paragraph{Contributions.}
\begin{enumerate}[leftmargin=*, itemsep=2pt]
  \item A backpropagation-free plastic prior on RSSM logits with norm-capped state $P_t$.
  \item A tri-timescale test-time stack: fast $P_t$, slow actor fine-tuning, and optional WM consolidation with coupled flush.
  \item Empirical characterization on gain shift: rapid trace growth and prior-error changes; control gains when fast plastic scaffolds slow actor learning; pilot evidence for $P_t$ drainage under consolidation.
\end{enumerate}


\section{Introduction}
\label{sec:introduction}

World models enable model-based agents to plan and act by predicting latent dynamics from experience~\cite{hafner2023mastering,hafner2020dream}.
Once trained, however, they are typically deployed with \emph{frozen} parameters.
In open-ended settings---wear on actuators, payload changes, or simulator-to-real gaps---the transition kernel internalized during pretraining ceases to match the environment the agent inhabits.
Prediction error rises, imagined rollouts become misleading, and control performance degrades.
\emph{Test-time adaptation} (TTA) seeks to recover competence after such shifts without returning to expensive offline training.

The dominant TTA recipe is gradient-based learning at deployment: collect post-shift trajectories, backpropagate a surrogate loss, and update network weights~\cite{nagabandi2018learning,zintgraf2021fast}.
This \emph{slow} pathway is expressive but incurs costs that matter in interactive control---batched updates, autograd through long horizons, and tuning to avoid catastrophic forgetting.
When the shift is sudden, the agent may act suboptimally for many steps before the first weight update.

Biological systems often separate a fast, local error-driven store from slower structural learning~\cite{mcclelland1995complementary}.
The fast store is not globally optimal; it provides coarse correction so behavior is not blind during transients.
We transpose this idea to world models: adaptation is not only a motor reflex, but a change in \emph{what the agent believes will happen next}.

We propose \textbf{Malleable World Models (MWM)}: frozen pretrained agents augmented with a \textbf{fast plastic prior residual} $P_t$, updated every step by a Hebbian rule without backprop through pretraining weights, plus optional \textbf{slow} actor and world-model updates from replay.
We evaluate on proprioceptive \texttt{walker-walk} under actuator gain shift ($1.0 \rightarrow 0.7 \rightarrow 1.0$).
Section~\ref{sec:method} details the method; Section~\ref{sec:experiments} reports detection, adaptation latency, and consolidation behavior.

\paragraph{Contributions.}
\begin{enumerate}[leftmargin=*, itemsep=2pt]
  \item A backpropagation-free plastic prior on RSSM logits with norm-capped state $P_t$.
  \item A tri-timescale test-time stack: fast $P_t$, slow actor fine-tuning, and optional WM consolidation with coupled flush.
  \item Empirical characterization on gain shift: rapid trace growth and prior-error changes; control gains when fast plastic scaffolds slow actor learning; pilot evidence for $P_t$ drainage under consolidation.
\end{enumerate}
, before % Method section for Malleable World Models (MWM)
% % Method section for Malleable World Models (MWM)
% % Method section for Malleable World Models (MWM)
% \input{method} from main.tex (requires amsmath, amssymb; optional booktabs)

\section{Method}
\label{sec:method}

We build on a pretrained model-based agent with a recurrent state-space model (RSSM)~\cite{hafner2020dream}.
At deployment, all pretrained parameters are \emph{frozen} for inference; test-time learning modifies either (i)~a separate fast plastic state~$P_t$, (ii)~the actor on a slow schedule, or (iii)~a subset of world-model weights on an even slower schedule.
Figure~\ref{fig:overview} (omitted here) summarizes the three pathways.
This section defines notation, the fast plastic layer, slow actor adaptation, consolidation with flush, and the gain-shift evaluation protocol.

%----------------------------------------------------------------------
\subsection{Preliminaries: frozen pretrained agent}
\label{sec:prelim}

Let $\theta$ denote the pretrained parameters of encoder $f_\theta$, RSSM $(h_\theta, g_\theta)$, decoder, and policy $\pi_\theta$, trained offline on task $\mathcal{T}_0$.
At test time we maintain frozen copies $(f, h, g, \pi) \leftarrow (f_\theta, h_\theta, g_\theta, \pi_\theta)$ used for action selection and latent rollout unless explicitly updated by a slow pathway (Section~\ref{sec:slow-actor}--\ref{sec:consolidation}).

Given observation $o_t$, the encoder produces embedding $e_t = f(o_t)$.
The RSSM maintains stochastic state $s_t$ and deterministic state $d_t$.
An observation step yields posterior logit $\ell^{\mathrm{post}}_t$ and updated $(s_t, d_t)$; the one-step prior over the \emph{next} stochastic state is
\begin{equation}
  \ell^{\mathrm{prior}}_t = g(d_t) \in \mathbb{R}^{K},
\end{equation}
where $K$ is the flattened stochastic logit dimension (for categorical latents, $K = N_{\mathrm{categ}} \times N_{\mathrm{classes}}$).
The policy samples $a_t \sim \pi(\cdot \mid \mathrm{feat}(s_t, d_t))$.

We evaluate under a controlled distribution shift that leaves observations and rewards in the same space but changes the environment's response to actions (Section~\ref{sec:shift}).

%----------------------------------------------------------------------
\subsection{Environment shift: actuator gain}
\label{sec:shift}

Physical simulator parameters (mass, friction, etc.) are held fixed.
Shift is induced by an \emph{action-gain} wrapper that scales the torque sent to the environment:
\begin{equation}
  \tilde{a}_t = \gamma_t \, a_t, \qquad \gamma_t \in (0, 1],
\end{equation}
where $a_t$ is the agent output and $\tilde{a}_t$ is executed.
This mimics weakened actuators or calibration drift without retraining the observation model.
Our default schedule uses three phases (Section~\ref{sec:protocol}): $\gamma=1.0$ (in-distribution), $\gamma=0.7$ (shift / adaptation), $\gamma=1.0$ (recovery).

%----------------------------------------------------------------------
\subsection{Fast pathway: plastic prior-logit residual}
\label{sec:fast-plastic}

\paragraph{State and readout.}
MWM augments the frozen prior with a test-time matrix $P_t \in \mathbb{R}^{K \times d}$ (one matrix per parallel environment), with no additional pretrained weights.
The plastic input is
\begin{equation}
  x_t = \big[\mathrm{flat}(s_t);\; d_{t-1};\; \hat{a}_{t-1}\big] \in \mathbb{R}^{d},
  \qquad
  d = K + |d| + |a|,
\end{equation}
where $\hat{a}_{t-1}$ is the previous action with each component divided by $\max(1, |a_{t-1,i}|)$ for numerical stability.
The corrected prior logit is
\begin{equation}
  \ell^{\mathrm{prior,pl}}_t = \ell^{\mathrm{prior}}_t + P_t x_t.
  \label{eq:plastic-read}
\end{equation}
During online interaction, Equation~\eqref{eq:plastic-read} affects the latent observation path only through the prior used to form prediction error; the policy still reads $(s_t, d_t)$ from the standard RSSM observation core with frozen weights.

\paragraph{Signed error.}
After the observation core produces $\ell^{\mathrm{post}}_t$, we define the clipped signed error
\begin{equation}
  e_t = \mathrm{clip}\!\left(\ell^{\mathrm{post}}_t - \ell^{\mathrm{prior,pl}}_t,\; \rho\right) \in \mathbb{R}^{K},
  \label{eq:signed-error}
\end{equation}
with per-vector norm clipped to $\rho$ (default $\rho=0.5$).
We track $\mathrm{MSE}_t = \|e_t\|_2^2 / K$ as a \emph{prior prediction error} diagnostic for shift detection.

\paragraph{Hebbian write (no backprop through $\theta$).}
$P_t$ is updated at \emph{every} environment step inside \texttt{torch.no\_grad()}:
\begin{align}
  \Delta P_t &= \eta \, g_t \, e_t x_t^\top - \lambda^{\mathrm{eff}}_t P_t,
  \label{eq:plastic-write}\\
  P_{t+1} &= \mathrm{Proj}_{\|\cdot\|_F \le \mu}\!\left(P_t + \Delta P_t\right),
  \label{eq:plastic-proj}
\end{align}
where $\eta$ is a plastic learning rate, $\mu$ is a Frobenius-norm cap (default $\mu=2$), and $g_t \in [0,1]$ is a surprise gate (below).
Equation~\eqref{eq:plastic-write} is a rank-one Hebbian term minus weight decay; it is local, causal, and does not backpropagate through the frozen backbone.

\paragraph{Effective decay and flush.}
The base decay is $\lambda_0$ (default $10^{-3}$).
When consolidation flush is enabled (Section~\ref{sec:consolidation}),
\begin{equation}
  \lambda^{\mathrm{eff}}_t = \lambda_0 + k_{\mathrm{flush}}\,\sigma\!\left(-\frac{\|e_t\|_2}{\tau}\right) + b_t,
  \label{eq:lambda-eff}
\end{equation}
where $\sigma$ is the logistic function, $\tau$ scales error magnitude (default $0.15$), $k_{\mathrm{flush}}$ is a config constant (default $0.5$), and $b_t \ge 0$ is a \emph{coupled flush boost} increased when slow world-model loss improves.
Large $\lambda^{\mathrm{eff}}_t$ accelerates draining of $P_t$ once the slow pathway has absorbed shift statistics.

\paragraph{Surprise gating.}
For gate mode \texttt{selective}, $g_t$ combines (i)~a normalized surprise statistic $z_t$ from an EMA of $\|e_t\|^2$ and (ii)~a direction-persistence score based on consecutive error alignment.
Modes \texttt{surprise\_only} and \texttt{always\_on} ablate parts of this logic; our main experiments use \texttt{always\_on} ($g_t \equiv 1$), so fast plasticity is continuously active and shift is reflected primarily in $e_t$ and $\|P_t\|_F$.

\paragraph{Episode reset.}
On \texttt{is\_first}, $P_t$ and gate statistics are zeroed for the corresponding batch index, but slow-updated actor / $\theta$ weights persist across episodes.

\paragraph{Capacity.}
For proprioceptive \texttt{walker-walk} with default RSSM ($32$ categorical factors $\times$ $16$ classes, $|d|=2048$, $|a|=6$), $P_t$ has shape $512 \times 2566$ (${\sim}1.3\times10^6$ floats, ${\sim}6\%$ of frozen parameter count).

%----------------------------------------------------------------------
\subsection{Online inference loop}
\label{sec:inference}

At each step with plastic inference enabled:
\begin{enumerate}[leftmargin=*, itemsep=1pt]
  \item Encode $e_t = f(o_t)$; run frozen RSSM observation core $\rightarrow (\ell^{\mathrm{post}}_t, s_t, d_t)$.
  \item Form $x_t$; read $P_t x_t$; compute $\ell^{\mathrm{prior,pl}}_t$ (Eq.~\eqref{eq:plastic-read}).
  \item Compute $e_t$ (Eq.~\eqref{eq:signed-error}); update $P_t$ (Eqs.~\eqref{eq:plastic-write}--\eqref{eq:plastic-proj}).
  \item Sample $a_t \sim \pi_\theta(\cdot \mid \mathrm{feat}(s_t, d_t))$ using the \emph{frozen} actor unless slow actor adaptation has modified $\pi_\theta$.
\end{enumerate}
Thus fast plasticity is a \emph{world-model correction} applied before the policy; it is not a direct motor reflex.
The agent's ``immediate'' response to shift is a change in latent beliefs, not an instantaneous change in policy weights.

%----------------------------------------------------------------------
\subsection{Slow pathway I: actor adaptation}
\label{sec:slow-actor}

Fast plasticity can realign one-step priors but does not directly optimize long-horizon return.
On phase $\mathcal{B}$ (shift), we collect transitions in a replay buffer $\mathcal{D}$ and periodically fine-tune only the actor via Dreamer-style imagination~\cite{hafner2020dream}.

\paragraph{Schedule.}
Every $\Delta_{\mathrm{adapt}}$ environment steps (default $\Delta_{\mathrm{adapt}}=50$), sample a sequence batch from $\mathcal{D}$ and perform $E$ gradient steps (default $E=4$) on the policy loss, with learning rate scaled by $\alpha_{\mathrm{actor}}$ (default $1.0$).

\paragraph{Imagination dynamics.}
Latent rollouts for the policy gradient use horizon $H$ (RSSM \texttt{imag\_horizon}).
Let \texttt{use\_plastic\_wm} control whether $P_t x$ is added during \emph{imagination}:
\begin{itemize}[leftmargin=*, itemsep=1pt]
  \item \textbf{Frozen-WM imagine} (\texttt{use\_plastic\_wm=false}, our default): rollouts follow $g$ without plastic read; reward/value/cont heads remain frozen copies.
  \item \textbf{Plastic-WM imagine} (\texttt{true}): a snapshot of $P_t$ is cloned and applied in each imagined step via Eq.~\eqref{eq:plastic-read}.
\end{itemize}
Online interaction always writes $P_t$ when plastic is enabled; the flag only affects actor gradients.
Empirically, frozen-WM imagine with online plastic in the replay distribution yields more stable control improvements than plastic-WM imagine.

\paragraph{Loss.}
Standard $\lambda$-return actor loss on imagined trajectories with frozen reward, value, and continuation heads; gradients are masked to actor parameters only (adaptive gradient clipping, AGC).

\paragraph{Online vs.\ offline plastic for actor learning.}
Condition~B enables online $P_t$ writes during $\mathcal{B}$; Condition~C (ablation) disables plastic inference but keeps actor adaptation, isolating whether live prior correction improves the data distribution for slow policy learning.

%----------------------------------------------------------------------
\subsection{Slow pathway II: world-model consolidation and flush}
\label{sec:consolidation}

When shifts persist, storing all compensatory structure in $P_t$ is bounded by $\mu$ and disappears on buffer reset.
The \emph{consolidation} pathway slowly fine-tunes world-model weights $\theta_{\mathrm{wm}} \subset \theta$ (encoder, RSSM, decoder / projector) on $\mathcal{D}$ while optionally increasing plastic decay to flush $P_t$ into $\theta_{\mathrm{wm}}$.

\paragraph{Slow WM update.}
Every $\Delta_{\mathrm{slow}}$ steps (default $100$), sample from $\mathcal{D}$ and minimize the pretraining world-model objective restricted to $\theta_{\mathrm{wm}}$:
dynamics KL, representation KL, and (for R2-Dreamer) Barlow Twins representation loss or (for Dreamer) reconstruction.
Learning rate is $\alpha_{\mathrm{wm}}$ times the pretraining base LR (default $\alpha_{\mathrm{wm}}=0.02$); only WM parameters receive gradients (AGC on WM parameters).

\paragraph{Frozen sync.}
After each slow update, live weights are copied into frozen inference modules:
$f, h, g \leftarrow f_\theta, h_\theta, g_\theta$.
The actor and reward heads are \emph{not} updated in this pathway unless separately adapted (Section~\ref{sec:slow-actor}).

\paragraph{Flush modes.}
\texttt{set\_consolidation\_mode} controls error-dependent decay (Eq.~\eqref{eq:lambda-eff}):
\begin{itemize}[leftmargin=*, itemsep=1pt]
  \item \texttt{off}: $k_{\mathrm{flush}}=0$, $b_t=0$.
  \item \texttt{coupled}: $k_{\mathrm{flush}}$ active; $b_t$ increases by $0.05$ when dyn loss improves vs.\ its EMA (cap $b_{\max}=0.5$), else decays by $0.98$.
  \item \texttt{error\_only}: $k_{\mathrm{flush}}$ active but $b_t \equiv 0$ (flush without slow $\theta$, a negative control).
\end{itemize}

\paragraph{Consolidation conditions (summary).}
\begin{table}[t]
  \centering
  \caption{Test-time conditions for consolidation experiments.}
  \label{tab:conditions}
  \begin{tabular}{@{}lcccc@{}}
    \toprule
    Condition & Online $P_t$ & Actor adapt & Slow $\theta_{\mathrm{wm}}$ & Flush mode \\
    \midrule
    A (baseline)     & \checkmark & \checkmark & --            & off \\
    B (slow WM only) & \checkmark & \checkmark & \checkmark    & off \\
    C (consolidate)  & \checkmark & \checkmark & \checkmark    & coupled \\
    D (trap)         & \checkmark & --         & --            & error\_only \\
    \bottomrule
  \end{tabular}
\end{table}

%----------------------------------------------------------------------
\subsection{Evaluation protocol}
\label{sec:protocol}

All test-time experiments share a three-phase rollout unless noted:
\begin{equation*}
  \mathcal{A}_1\ (\gamma{=}1.0,\ T_1{=}200)
  \;\rightarrow\;
  \mathcal{B}\ (\gamma{=}0.7,\ T_2{=}1000)
  \;\rightarrow\;
  \mathcal{A}_2\ (\gamma{=}1.0,\ T_3{=}300).
\end{equation*}
$\mathcal{A}_1$ establishes pre-shift performance; $\mathcal{B}$ is the adaptation window (actor / slow WM updates occur only here by default); $\mathcal{A}_2$ tests recovery when gain returns to nominal.

\paragraph{Metrics.}
\begin{itemize}[leftmargin=*, itemsep=1pt]
  \item \textbf{Return}: mean per-step reward in each phase.
  \item \textbf{Prior MSE}: $\mathrm{MSE}_t$ from Eq.~\eqref{eq:signed-error}, averaged over steps (WM alignment).
  \item \textbf{Trace norm}: $\|P_t\|_F$, capped at $\mu$ (fast-store magnitude).
  \item \textbf{Recovery time}: smallest $t \in \mathcal{B}$ such that the rolling mean return exceeds $0.9 \times$ the $\mathcal{A}_1$ baseline (reported per seed).
  \item \textbf{Slow dynamics loss}: RSSM KL during $\theta_{\mathrm{wm}}$ updates (consolidation diagnostics).
\end{itemize}

\paragraph{Implementation.}
Built on R2-Dreamer with test-time modules \texttt{LogitPlasticResidual}, \texttt{plastic\_actor\_adapt.py}, and \texttt{plastic\_consolidation\_eval.py}.
Pretrained checkpoint: proprioceptive \texttt{walker-walk}, frozen at eval.
Plastic defaults: $\eta=0.05$, $\lambda_0=10^{-3}$, $\mu=2$, \texttt{gate\_mode=always\_on}.
Actor-adapt default (\texttt{cfg0}): $\Delta_{\mathrm{adapt}}=50$, $E=4$, $\alpha_{\mathrm{actor}}=1$, frozen-WM imagine.

%----------------------------------------------------------------------
\subsection{Design rationale}
\label{sec:rationale}

Equation~\eqref{eq:plastic-write} implements a \emph{fast, imprecise} store: it reacts on the first post-shift errors with $O(Kd)$ state and no autograd.
Actor and WM gradient updates are \emph{slow, selective} stores: they require batches, imagination, and careful learning-rate / AGC tuning but can improve long-horizon control and, in principle, absorb persistent structure into $\theta_{\mathrm{wm}}$.
Coupled flush (Eq.~\eqref{eq:lambda-eff}) expresses the hypothesis that fast compensation should retire once slow parameters explain the shifted dynamics---a computational analogue of complementary learning systems~\cite{mcclelland1995complementary}.
Which parts of this stack are necessary for detectable, fast, and lasting adaptation is an empirical question addressed in Section~\ref{sec:experiments}.
 from main.tex (requires amsmath, amssymb; optional booktabs)

\section{Method}
\label{sec:method}

We build on a pretrained model-based agent with a recurrent state-space model (RSSM)~\cite{hafner2020dream}.
At deployment, all pretrained parameters are \emph{frozen} for inference; test-time learning modifies either (i)~a separate fast plastic state~$P_t$, (ii)~the actor on a slow schedule, or (iii)~a subset of world-model weights on an even slower schedule.
Figure~\ref{fig:overview} (omitted here) summarizes the three pathways.
This section defines notation, the fast plastic layer, slow actor adaptation, consolidation with flush, and the gain-shift evaluation protocol.

%----------------------------------------------------------------------
\subsection{Preliminaries: frozen pretrained agent}
\label{sec:prelim}

Let $\theta$ denote the pretrained parameters of encoder $f_\theta$, RSSM $(h_\theta, g_\theta)$, decoder, and policy $\pi_\theta$, trained offline on task $\mathcal{T}_0$.
At test time we maintain frozen copies $(f, h, g, \pi) \leftarrow (f_\theta, h_\theta, g_\theta, \pi_\theta)$ used for action selection and latent rollout unless explicitly updated by a slow pathway (Section~\ref{sec:slow-actor}--\ref{sec:consolidation}).

Given observation $o_t$, the encoder produces embedding $e_t = f(o_t)$.
The RSSM maintains stochastic state $s_t$ and deterministic state $d_t$.
An observation step yields posterior logit $\ell^{\mathrm{post}}_t$ and updated $(s_t, d_t)$; the one-step prior over the \emph{next} stochastic state is
\begin{equation}
  \ell^{\mathrm{prior}}_t = g(d_t) \in \mathbb{R}^{K},
\end{equation}
where $K$ is the flattened stochastic logit dimension (for categorical latents, $K = N_{\mathrm{categ}} \times N_{\mathrm{classes}}$).
The policy samples $a_t \sim \pi(\cdot \mid \mathrm{feat}(s_t, d_t))$.

We evaluate under a controlled distribution shift that leaves observations and rewards in the same space but changes the environment's response to actions (Section~\ref{sec:shift}).

%----------------------------------------------------------------------
\subsection{Environment shift: actuator gain}
\label{sec:shift}

Physical simulator parameters (mass, friction, etc.) are held fixed.
Shift is induced by an \emph{action-gain} wrapper that scales the torque sent to the environment:
\begin{equation}
  \tilde{a}_t = \gamma_t \, a_t, \qquad \gamma_t \in (0, 1],
\end{equation}
where $a_t$ is the agent output and $\tilde{a}_t$ is executed.
This mimics weakened actuators or calibration drift without retraining the observation model.
Our default schedule uses three phases (Section~\ref{sec:protocol}): $\gamma=1.0$ (in-distribution), $\gamma=0.7$ (shift / adaptation), $\gamma=1.0$ (recovery).

%----------------------------------------------------------------------
\subsection{Fast pathway: plastic prior-logit residual}
\label{sec:fast-plastic}

\paragraph{State and readout.}
MWM augments the frozen prior with a test-time matrix $P_t \in \mathbb{R}^{K \times d}$ (one matrix per parallel environment), with no additional pretrained weights.
The plastic input is
\begin{equation}
  x_t = \big[\mathrm{flat}(s_t);\; d_{t-1};\; \hat{a}_{t-1}\big] \in \mathbb{R}^{d},
  \qquad
  d = K + |d| + |a|,
\end{equation}
where $\hat{a}_{t-1}$ is the previous action with each component divided by $\max(1, |a_{t-1,i}|)$ for numerical stability.
The corrected prior logit is
\begin{equation}
  \ell^{\mathrm{prior,pl}}_t = \ell^{\mathrm{prior}}_t + P_t x_t.
  \label{eq:plastic-read}
\end{equation}
During online interaction, Equation~\eqref{eq:plastic-read} affects the latent observation path only through the prior used to form prediction error; the policy still reads $(s_t, d_t)$ from the standard RSSM observation core with frozen weights.

\paragraph{Signed error.}
After the observation core produces $\ell^{\mathrm{post}}_t$, we define the clipped signed error
\begin{equation}
  e_t = \mathrm{clip}\!\left(\ell^{\mathrm{post}}_t - \ell^{\mathrm{prior,pl}}_t,\; \rho\right) \in \mathbb{R}^{K},
  \label{eq:signed-error}
\end{equation}
with per-vector norm clipped to $\rho$ (default $\rho=0.5$).
We track $\mathrm{MSE}_t = \|e_t\|_2^2 / K$ as a \emph{prior prediction error} diagnostic for shift detection.

\paragraph{Hebbian write (no backprop through $\theta$).}
$P_t$ is updated at \emph{every} environment step inside \texttt{torch.no\_grad()}:
\begin{align}
  \Delta P_t &= \eta \, g_t \, e_t x_t^\top - \lambda^{\mathrm{eff}}_t P_t,
  \label{eq:plastic-write}\\
  P_{t+1} &= \mathrm{Proj}_{\|\cdot\|_F \le \mu}\!\left(P_t + \Delta P_t\right),
  \label{eq:plastic-proj}
\end{align}
where $\eta$ is a plastic learning rate, $\mu$ is a Frobenius-norm cap (default $\mu=2$), and $g_t \in [0,1]$ is a surprise gate (below).
Equation~\eqref{eq:plastic-write} is a rank-one Hebbian term minus weight decay; it is local, causal, and does not backpropagate through the frozen backbone.

\paragraph{Effective decay and flush.}
The base decay is $\lambda_0$ (default $10^{-3}$).
When consolidation flush is enabled (Section~\ref{sec:consolidation}),
\begin{equation}
  \lambda^{\mathrm{eff}}_t = \lambda_0 + k_{\mathrm{flush}}\,\sigma\!\left(-\frac{\|e_t\|_2}{\tau}\right) + b_t,
  \label{eq:lambda-eff}
\end{equation}
where $\sigma$ is the logistic function, $\tau$ scales error magnitude (default $0.15$), $k_{\mathrm{flush}}$ is a config constant (default $0.5$), and $b_t \ge 0$ is a \emph{coupled flush boost} increased when slow world-model loss improves.
Large $\lambda^{\mathrm{eff}}_t$ accelerates draining of $P_t$ once the slow pathway has absorbed shift statistics.

\paragraph{Surprise gating.}
For gate mode \texttt{selective}, $g_t$ combines (i)~a normalized surprise statistic $z_t$ from an EMA of $\|e_t\|^2$ and (ii)~a direction-persistence score based on consecutive error alignment.
Modes \texttt{surprise\_only} and \texttt{always\_on} ablate parts of this logic; our main experiments use \texttt{always\_on} ($g_t \equiv 1$), so fast plasticity is continuously active and shift is reflected primarily in $e_t$ and $\|P_t\|_F$.

\paragraph{Episode reset.}
On \texttt{is\_first}, $P_t$ and gate statistics are zeroed for the corresponding batch index, but slow-updated actor / $\theta$ weights persist across episodes.

\paragraph{Capacity.}
For proprioceptive \texttt{walker-walk} with default RSSM ($32$ categorical factors $\times$ $16$ classes, $|d|=2048$, $|a|=6$), $P_t$ has shape $512 \times 2566$ (${\sim}1.3\times10^6$ floats, ${\sim}6\%$ of frozen parameter count).

%----------------------------------------------------------------------
\subsection{Online inference loop}
\label{sec:inference}

At each step with plastic inference enabled:
\begin{enumerate}[leftmargin=*, itemsep=1pt]
  \item Encode $e_t = f(o_t)$; run frozen RSSM observation core $\rightarrow (\ell^{\mathrm{post}}_t, s_t, d_t)$.
  \item Form $x_t$; read $P_t x_t$; compute $\ell^{\mathrm{prior,pl}}_t$ (Eq.~\eqref{eq:plastic-read}).
  \item Compute $e_t$ (Eq.~\eqref{eq:signed-error}); update $P_t$ (Eqs.~\eqref{eq:plastic-write}--\eqref{eq:plastic-proj}).
  \item Sample $a_t \sim \pi_\theta(\cdot \mid \mathrm{feat}(s_t, d_t))$ using the \emph{frozen} actor unless slow actor adaptation has modified $\pi_\theta$.
\end{enumerate}
Thus fast plasticity is a \emph{world-model correction} applied before the policy; it is not a direct motor reflex.
The agent's ``immediate'' response to shift is a change in latent beliefs, not an instantaneous change in policy weights.

%----------------------------------------------------------------------
\subsection{Slow pathway I: actor adaptation}
\label{sec:slow-actor}

Fast plasticity can realign one-step priors but does not directly optimize long-horizon return.
On phase $\mathcal{B}$ (shift), we collect transitions in a replay buffer $\mathcal{D}$ and periodically fine-tune only the actor via Dreamer-style imagination~\cite{hafner2020dream}.

\paragraph{Schedule.}
Every $\Delta_{\mathrm{adapt}}$ environment steps (default $\Delta_{\mathrm{adapt}}=50$), sample a sequence batch from $\mathcal{D}$ and perform $E$ gradient steps (default $E=4$) on the policy loss, with learning rate scaled by $\alpha_{\mathrm{actor}}$ (default $1.0$).

\paragraph{Imagination dynamics.}
Latent rollouts for the policy gradient use horizon $H$ (RSSM \texttt{imag\_horizon}).
Let \texttt{use\_plastic\_wm} control whether $P_t x$ is added during \emph{imagination}:
\begin{itemize}[leftmargin=*, itemsep=1pt]
  \item \textbf{Frozen-WM imagine} (\texttt{use\_plastic\_wm=false}, our default): rollouts follow $g$ without plastic read; reward/value/cont heads remain frozen copies.
  \item \textbf{Plastic-WM imagine} (\texttt{true}): a snapshot of $P_t$ is cloned and applied in each imagined step via Eq.~\eqref{eq:plastic-read}.
\end{itemize}
Online interaction always writes $P_t$ when plastic is enabled; the flag only affects actor gradients.
Empirically, frozen-WM imagine with online plastic in the replay distribution yields more stable control improvements than plastic-WM imagine.

\paragraph{Loss.}
Standard $\lambda$-return actor loss on imagined trajectories with frozen reward, value, and continuation heads; gradients are masked to actor parameters only (adaptive gradient clipping, AGC).

\paragraph{Online vs.\ offline plastic for actor learning.}
Condition~B enables online $P_t$ writes during $\mathcal{B}$; Condition~C (ablation) disables plastic inference but keeps actor adaptation, isolating whether live prior correction improves the data distribution for slow policy learning.

%----------------------------------------------------------------------
\subsection{Slow pathway II: world-model consolidation and flush}
\label{sec:consolidation}

When shifts persist, storing all compensatory structure in $P_t$ is bounded by $\mu$ and disappears on buffer reset.
The \emph{consolidation} pathway slowly fine-tunes world-model weights $\theta_{\mathrm{wm}} \subset \theta$ (encoder, RSSM, decoder / projector) on $\mathcal{D}$ while optionally increasing plastic decay to flush $P_t$ into $\theta_{\mathrm{wm}}$.

\paragraph{Slow WM update.}
Every $\Delta_{\mathrm{slow}}$ steps (default $100$), sample from $\mathcal{D}$ and minimize the pretraining world-model objective restricted to $\theta_{\mathrm{wm}}$:
dynamics KL, representation KL, and (for R2-Dreamer) Barlow Twins representation loss or (for Dreamer) reconstruction.
Learning rate is $\alpha_{\mathrm{wm}}$ times the pretraining base LR (default $\alpha_{\mathrm{wm}}=0.02$); only WM parameters receive gradients (AGC on WM parameters).

\paragraph{Frozen sync.}
After each slow update, live weights are copied into frozen inference modules:
$f, h, g \leftarrow f_\theta, h_\theta, g_\theta$.
The actor and reward heads are \emph{not} updated in this pathway unless separately adapted (Section~\ref{sec:slow-actor}).

\paragraph{Flush modes.}
\texttt{set\_consolidation\_mode} controls error-dependent decay (Eq.~\eqref{eq:lambda-eff}):
\begin{itemize}[leftmargin=*, itemsep=1pt]
  \item \texttt{off}: $k_{\mathrm{flush}}=0$, $b_t=0$.
  \item \texttt{coupled}: $k_{\mathrm{flush}}$ active; $b_t$ increases by $0.05$ when dyn loss improves vs.\ its EMA (cap $b_{\max}=0.5$), else decays by $0.98$.
  \item \texttt{error\_only}: $k_{\mathrm{flush}}$ active but $b_t \equiv 0$ (flush without slow $\theta$, a negative control).
\end{itemize}

\paragraph{Consolidation conditions (summary).}
\begin{table}[t]
  \centering
  \caption{Test-time conditions for consolidation experiments.}
  \label{tab:conditions}
  \begin{tabular}{@{}lcccc@{}}
    \toprule
    Condition & Online $P_t$ & Actor adapt & Slow $\theta_{\mathrm{wm}}$ & Flush mode \\
    \midrule
    A (baseline)     & \checkmark & \checkmark & --            & off \\
    B (slow WM only) & \checkmark & \checkmark & \checkmark    & off \\
    C (consolidate)  & \checkmark & \checkmark & \checkmark    & coupled \\
    D (trap)         & \checkmark & --         & --            & error\_only \\
    \bottomrule
  \end{tabular}
\end{table}

%----------------------------------------------------------------------
\subsection{Evaluation protocol}
\label{sec:protocol}

All test-time experiments share a three-phase rollout unless noted:
\begin{equation*}
  \mathcal{A}_1\ (\gamma{=}1.0,\ T_1{=}200)
  \;\rightarrow\;
  \mathcal{B}\ (\gamma{=}0.7,\ T_2{=}1000)
  \;\rightarrow\;
  \mathcal{A}_2\ (\gamma{=}1.0,\ T_3{=}300).
\end{equation*}
$\mathcal{A}_1$ establishes pre-shift performance; $\mathcal{B}$ is the adaptation window (actor / slow WM updates occur only here by default); $\mathcal{A}_2$ tests recovery when gain returns to nominal.

\paragraph{Metrics.}
\begin{itemize}[leftmargin=*, itemsep=1pt]
  \item \textbf{Return}: mean per-step reward in each phase.
  \item \textbf{Prior MSE}: $\mathrm{MSE}_t$ from Eq.~\eqref{eq:signed-error}, averaged over steps (WM alignment).
  \item \textbf{Trace norm}: $\|P_t\|_F$, capped at $\mu$ (fast-store magnitude).
  \item \textbf{Recovery time}: smallest $t \in \mathcal{B}$ such that the rolling mean return exceeds $0.9 \times$ the $\mathcal{A}_1$ baseline (reported per seed).
  \item \textbf{Slow dynamics loss}: RSSM KL during $\theta_{\mathrm{wm}}$ updates (consolidation diagnostics).
\end{itemize}

\paragraph{Implementation.}
Built on R2-Dreamer with test-time modules \texttt{LogitPlasticResidual}, \texttt{plastic\_actor\_adapt.py}, and \texttt{plastic\_consolidation\_eval.py}.
Pretrained checkpoint: proprioceptive \texttt{walker-walk}, frozen at eval.
Plastic defaults: $\eta=0.05$, $\lambda_0=10^{-3}$, $\mu=2$, \texttt{gate\_mode=always\_on}.
Actor-adapt default (\texttt{cfg0}): $\Delta_{\mathrm{adapt}}=50$, $E=4$, $\alpha_{\mathrm{actor}}=1$, frozen-WM imagine.

%----------------------------------------------------------------------
\subsection{Design rationale}
\label{sec:rationale}

Equation~\eqref{eq:plastic-write} implements a \emph{fast, imprecise} store: it reacts on the first post-shift errors with $O(Kd)$ state and no autograd.
Actor and WM gradient updates are \emph{slow, selective} stores: they require batches, imagination, and careful learning-rate / AGC tuning but can improve long-horizon control and, in principle, absorb persistent structure into $\theta_{\mathrm{wm}}$.
Coupled flush (Eq.~\eqref{eq:lambda-eff}) expresses the hypothesis that fast compensation should retire once slow parameters explain the shifted dynamics---a computational analogue of complementary learning systems~\cite{mcclelland1995complementary}.
Which parts of this stack are necessary for detectable, fast, and lasting adaptation is an empirical question addressed in Section~\ref{sec:experiments}.
 from main.tex (requires amsmath, amssymb; optional booktabs)

\section{Method}
\label{sec:method}

We build on a pretrained model-based agent with a recurrent state-space model (RSSM)~\cite{hafner2020dream}.
At deployment, all pretrained parameters are \emph{frozen} for inference; test-time learning modifies either (i)~a separate fast plastic state~$P_t$, (ii)~the actor on a slow schedule, or (iii)~a subset of world-model weights on an even slower schedule.
Figure~\ref{fig:overview} (omitted here) summarizes the three pathways.
This section defines notation, the fast plastic layer, slow actor adaptation, consolidation with flush, and the gain-shift evaluation protocol.

%----------------------------------------------------------------------
\subsection{Preliminaries: frozen pretrained agent}
\label{sec:prelim}

Let $\theta$ denote the pretrained parameters of encoder $f_\theta$, RSSM $(h_\theta, g_\theta)$, decoder, and policy $\pi_\theta$, trained offline on task $\mathcal{T}_0$.
At test time we maintain frozen copies $(f, h, g, \pi) \leftarrow (f_\theta, h_\theta, g_\theta, \pi_\theta)$ used for action selection and latent rollout unless explicitly updated by a slow pathway (Section~\ref{sec:slow-actor}--\ref{sec:consolidation}).

Given observation $o_t$, the encoder produces embedding $e_t = f(o_t)$.
The RSSM maintains stochastic state $s_t$ and deterministic state $d_t$.
An observation step yields posterior logit $\ell^{\mathrm{post}}_t$ and updated $(s_t, d_t)$; the one-step prior over the \emph{next} stochastic state is
\begin{equation}
  \ell^{\mathrm{prior}}_t = g(d_t) \in \mathbb{R}^{K},
\end{equation}
where $K$ is the flattened stochastic logit dimension (for categorical latents, $K = N_{\mathrm{categ}} \times N_{\mathrm{classes}}$).
The policy samples $a_t \sim \pi(\cdot \mid \mathrm{feat}(s_t, d_t))$.

We evaluate under a controlled distribution shift that leaves observations and rewards in the same space but changes the environment's response to actions (Section~\ref{sec:shift}).

%----------------------------------------------------------------------
\subsection{Environment shift: actuator gain}
\label{sec:shift}

Physical simulator parameters (mass, friction, etc.) are held fixed.
Shift is induced by an \emph{action-gain} wrapper that scales the torque sent to the environment:
\begin{equation}
  \tilde{a}_t = \gamma_t \, a_t, \qquad \gamma_t \in (0, 1],
\end{equation}
where $a_t$ is the agent output and $\tilde{a}_t$ is executed.
This mimics weakened actuators or calibration drift without retraining the observation model.
Our default schedule uses three phases (Section~\ref{sec:protocol}): $\gamma=1.0$ (in-distribution), $\gamma=0.7$ (shift / adaptation), $\gamma=1.0$ (recovery).

%----------------------------------------------------------------------
\subsection{Fast pathway: plastic prior-logit residual}
\label{sec:fast-plastic}

\paragraph{State and readout.}
MWM augments the frozen prior with a test-time matrix $P_t \in \mathbb{R}^{K \times d}$ (one matrix per parallel environment), with no additional pretrained weights.
The plastic input is
\begin{equation}
  x_t = \big[\mathrm{flat}(s_t);\; d_{t-1};\; \hat{a}_{t-1}\big] \in \mathbb{R}^{d},
  \qquad
  d = K + |d| + |a|,
\end{equation}
where $\hat{a}_{t-1}$ is the previous action with each component divided by $\max(1, |a_{t-1,i}|)$ for numerical stability.
The corrected prior logit is
\begin{equation}
  \ell^{\mathrm{prior,pl}}_t = \ell^{\mathrm{prior}}_t + P_t x_t.
  \label{eq:plastic-read}
\end{equation}
During online interaction, Equation~\eqref{eq:plastic-read} affects the latent observation path only through the prior used to form prediction error; the policy still reads $(s_t, d_t)$ from the standard RSSM observation core with frozen weights.

\paragraph{Signed error.}
After the observation core produces $\ell^{\mathrm{post}}_t$, we define the clipped signed error
\begin{equation}
  e_t = \mathrm{clip}\!\left(\ell^{\mathrm{post}}_t - \ell^{\mathrm{prior,pl}}_t,\; \rho\right) \in \mathbb{R}^{K},
  \label{eq:signed-error}
\end{equation}
with per-vector norm clipped to $\rho$ (default $\rho=0.5$).
We track $\mathrm{MSE}_t = \|e_t\|_2^2 / K$ as a \emph{prior prediction error} diagnostic for shift detection.

\paragraph{Hebbian write (no backprop through $\theta$).}
$P_t$ is updated at \emph{every} environment step inside \texttt{torch.no\_grad()}:
\begin{align}
  \Delta P_t &= \eta \, g_t \, e_t x_t^\top - \lambda^{\mathrm{eff}}_t P_t,
  \label{eq:plastic-write}\\
  P_{t+1} &= \mathrm{Proj}_{\|\cdot\|_F \le \mu}\!\left(P_t + \Delta P_t\right),
  \label{eq:plastic-proj}
\end{align}
where $\eta$ is a plastic learning rate, $\mu$ is a Frobenius-norm cap (default $\mu=2$), and $g_t \in [0,1]$ is a surprise gate (below).
Equation~\eqref{eq:plastic-write} is a rank-one Hebbian term minus weight decay; it is local, causal, and does not backpropagate through the frozen backbone.

\paragraph{Effective decay and flush.}
The base decay is $\lambda_0$ (default $10^{-3}$).
When consolidation flush is enabled (Section~\ref{sec:consolidation}),
\begin{equation}
  \lambda^{\mathrm{eff}}_t = \lambda_0 + k_{\mathrm{flush}}\,\sigma\!\left(-\frac{\|e_t\|_2}{\tau}\right) + b_t,
  \label{eq:lambda-eff}
\end{equation}
where $\sigma$ is the logistic function, $\tau$ scales error magnitude (default $0.15$), $k_{\mathrm{flush}}$ is a config constant (default $0.5$), and $b_t \ge 0$ is a \emph{coupled flush boost} increased when slow world-model loss improves.
Large $\lambda^{\mathrm{eff}}_t$ accelerates draining of $P_t$ once the slow pathway has absorbed shift statistics.

\paragraph{Surprise gating.}
For gate mode \texttt{selective}, $g_t$ combines (i)~a normalized surprise statistic $z_t$ from an EMA of $\|e_t\|^2$ and (ii)~a direction-persistence score based on consecutive error alignment.
Modes \texttt{surprise\_only} and \texttt{always\_on} ablate parts of this logic; our main experiments use \texttt{always\_on} ($g_t \equiv 1$), so fast plasticity is continuously active and shift is reflected primarily in $e_t$ and $\|P_t\|_F$.

\paragraph{Episode reset.}
On \texttt{is\_first}, $P_t$ and gate statistics are zeroed for the corresponding batch index, but slow-updated actor / $\theta$ weights persist across episodes.

\paragraph{Capacity.}
For proprioceptive \texttt{walker-walk} with default RSSM ($32$ categorical factors $\times$ $16$ classes, $|d|=2048$, $|a|=6$), $P_t$ has shape $512 \times 2566$ (${\sim}1.3\times10^6$ floats, ${\sim}6\%$ of frozen parameter count).

%----------------------------------------------------------------------
\subsection{Online inference loop}
\label{sec:inference}

At each step with plastic inference enabled:
\begin{enumerate}[leftmargin=*, itemsep=1pt]
  \item Encode $e_t = f(o_t)$; run frozen RSSM observation core $\rightarrow (\ell^{\mathrm{post}}_t, s_t, d_t)$.
  \item Form $x_t$; read $P_t x_t$; compute $\ell^{\mathrm{prior,pl}}_t$ (Eq.~\eqref{eq:plastic-read}).
  \item Compute $e_t$ (Eq.~\eqref{eq:signed-error}); update $P_t$ (Eqs.~\eqref{eq:plastic-write}--\eqref{eq:plastic-proj}).
  \item Sample $a_t \sim \pi_\theta(\cdot \mid \mathrm{feat}(s_t, d_t))$ using the \emph{frozen} actor unless slow actor adaptation has modified $\pi_\theta$.
\end{enumerate}
Thus fast plasticity is a \emph{world-model correction} applied before the policy; it is not a direct motor reflex.
The agent's ``immediate'' response to shift is a change in latent beliefs, not an instantaneous change in policy weights.

%----------------------------------------------------------------------
\subsection{Slow pathway I: actor adaptation}
\label{sec:slow-actor}

Fast plasticity can realign one-step priors but does not directly optimize long-horizon return.
On phase $\mathcal{B}$ (shift), we collect transitions in a replay buffer $\mathcal{D}$ and periodically fine-tune only the actor via Dreamer-style imagination~\cite{hafner2020dream}.

\paragraph{Schedule.}
Every $\Delta_{\mathrm{adapt}}$ environment steps (default $\Delta_{\mathrm{adapt}}=50$), sample a sequence batch from $\mathcal{D}$ and perform $E$ gradient steps (default $E=4$) on the policy loss, with learning rate scaled by $\alpha_{\mathrm{actor}}$ (default $1.0$).

\paragraph{Imagination dynamics.}
Latent rollouts for the policy gradient use horizon $H$ (RSSM \texttt{imag\_horizon}).
Let \texttt{use\_plastic\_wm} control whether $P_t x$ is added during \emph{imagination}:
\begin{itemize}[leftmargin=*, itemsep=1pt]
  \item \textbf{Frozen-WM imagine} (\texttt{use\_plastic\_wm=false}, our default): rollouts follow $g$ without plastic read; reward/value/cont heads remain frozen copies.
  \item \textbf{Plastic-WM imagine} (\texttt{true}): a snapshot of $P_t$ is cloned and applied in each imagined step via Eq.~\eqref{eq:plastic-read}.
\end{itemize}
Online interaction always writes $P_t$ when plastic is enabled; the flag only affects actor gradients.
Empirically, frozen-WM imagine with online plastic in the replay distribution yields more stable control improvements than plastic-WM imagine.

\paragraph{Loss.}
Standard $\lambda$-return actor loss on imagined trajectories with frozen reward, value, and continuation heads; gradients are masked to actor parameters only (adaptive gradient clipping, AGC).

\paragraph{Online vs.\ offline plastic for actor learning.}
Condition~B enables online $P_t$ writes during $\mathcal{B}$; Condition~C (ablation) disables plastic inference but keeps actor adaptation, isolating whether live prior correction improves the data distribution for slow policy learning.

%----------------------------------------------------------------------
\subsection{Slow pathway II: world-model consolidation and flush}
\label{sec:consolidation}

When shifts persist, storing all compensatory structure in $P_t$ is bounded by $\mu$ and disappears on buffer reset.
The \emph{consolidation} pathway slowly fine-tunes world-model weights $\theta_{\mathrm{wm}} \subset \theta$ (encoder, RSSM, decoder / projector) on $\mathcal{D}$ while optionally increasing plastic decay to flush $P_t$ into $\theta_{\mathrm{wm}}$.

\paragraph{Slow WM update.}
Every $\Delta_{\mathrm{slow}}$ steps (default $100$), sample from $\mathcal{D}$ and minimize the pretraining world-model objective restricted to $\theta_{\mathrm{wm}}$:
dynamics KL, representation KL, and (for R2-Dreamer) Barlow Twins representation loss or (for Dreamer) reconstruction.
Learning rate is $\alpha_{\mathrm{wm}}$ times the pretraining base LR (default $\alpha_{\mathrm{wm}}=0.02$); only WM parameters receive gradients (AGC on WM parameters).

\paragraph{Frozen sync.}
After each slow update, live weights are copied into frozen inference modules:
$f, h, g \leftarrow f_\theta, h_\theta, g_\theta$.
The actor and reward heads are \emph{not} updated in this pathway unless separately adapted (Section~\ref{sec:slow-actor}).

\paragraph{Flush modes.}
\texttt{set\_consolidation\_mode} controls error-dependent decay (Eq.~\eqref{eq:lambda-eff}):
\begin{itemize}[leftmargin=*, itemsep=1pt]
  \item \texttt{off}: $k_{\mathrm{flush}}=0$, $b_t=0$.
  \item \texttt{coupled}: $k_{\mathrm{flush}}$ active; $b_t$ increases by $0.05$ when dyn loss improves vs.\ its EMA (cap $b_{\max}=0.5$), else decays by $0.98$.
  \item \texttt{error\_only}: $k_{\mathrm{flush}}$ active but $b_t \equiv 0$ (flush without slow $\theta$, a negative control).
\end{itemize}

\paragraph{Consolidation conditions (summary).}
\begin{table}[t]
  \centering
  \caption{Test-time conditions for consolidation experiments.}
  \label{tab:conditions}
  \begin{tabular}{@{}lcccc@{}}
    \toprule
    Condition & Online $P_t$ & Actor adapt & Slow $\theta_{\mathrm{wm}}$ & Flush mode \\
    \midrule
    A (baseline)     & \checkmark & \checkmark & --            & off \\
    B (slow WM only) & \checkmark & \checkmark & \checkmark    & off \\
    C (consolidate)  & \checkmark & \checkmark & \checkmark    & coupled \\
    D (trap)         & \checkmark & --         & --            & error\_only \\
    \bottomrule
  \end{tabular}
\end{table}

%----------------------------------------------------------------------
\subsection{Evaluation protocol}
\label{sec:protocol}

All test-time experiments share a three-phase rollout unless noted:
\begin{equation*}
  \mathcal{A}_1\ (\gamma{=}1.0,\ T_1{=}200)
  \;\rightarrow\;
  \mathcal{B}\ (\gamma{=}0.7,\ T_2{=}1000)
  \;\rightarrow\;
  \mathcal{A}_2\ (\gamma{=}1.0,\ T_3{=}300).
\end{equation*}
$\mathcal{A}_1$ establishes pre-shift performance; $\mathcal{B}$ is the adaptation window (actor / slow WM updates occur only here by default); $\mathcal{A}_2$ tests recovery when gain returns to nominal.

\paragraph{Metrics.}
\begin{itemize}[leftmargin=*, itemsep=1pt]
  \item \textbf{Return}: mean per-step reward in each phase.
  \item \textbf{Prior MSE}: $\mathrm{MSE}_t$ from Eq.~\eqref{eq:signed-error}, averaged over steps (WM alignment).
  \item \textbf{Trace norm}: $\|P_t\|_F$, capped at $\mu$ (fast-store magnitude).
  \item \textbf{Recovery time}: smallest $t \in \mathcal{B}$ such that the rolling mean return exceeds $0.9 \times$ the $\mathcal{A}_1$ baseline (reported per seed).
  \item \textbf{Slow dynamics loss}: RSSM KL during $\theta_{\mathrm{wm}}$ updates (consolidation diagnostics).
\end{itemize}

\paragraph{Implementation.}
Built on R2-Dreamer with test-time modules \texttt{LogitPlasticResidual}, \texttt{plastic\_actor\_adapt.py}, and \texttt{plastic\_consolidation\_eval.py}.
Pretrained checkpoint: proprioceptive \texttt{walker-walk}, frozen at eval.
Plastic defaults: $\eta=0.05$, $\lambda_0=10^{-3}$, $\mu=2$, \texttt{gate\_mode=always\_on}.
Actor-adapt default (\texttt{cfg0}): $\Delta_{\mathrm{adapt}}=50$, $E=4$, $\alpha_{\mathrm{actor}}=1$, frozen-WM imagine.

%----------------------------------------------------------------------
\subsection{Design rationale}
\label{sec:rationale}

Equation~\eqref{eq:plastic-write} implements a \emph{fast, imprecise} store: it reacts on the first post-shift errors with $O(Kd)$ state and no autograd.
Actor and WM gradient updates are \emph{slow, selective} stores: they require batches, imagination, and careful learning-rate / AGC tuning but can improve long-horizon control and, in principle, absorb persistent structure into $\theta_{\mathrm{wm}}$.
Coupled flush (Eq.~\eqref{eq:lambda-eff}) expresses the hypothesis that fast compensation should retire once slow parameters explain the shifted dynamics---a computational analogue of complementary learning systems~\cite{mcclelland1995complementary}.
Which parts of this stack are necessary for detectable, fast, and lasting adaptation is an empirical question addressed in Section~\ref{sec:experiments}.


\section{Related Work}
\label{sec:related}

MWM sits at the intersection of model-based reinforcement learning, test-time adaptation, and dual-timescale memory.
We review each line of work and clarify what our fast plastic prior adds.

%----------------------------------------------------------------------
\subsection{World models and model-based RL}
\label{sec:related-wm}

Learned dynamics models enable sample-efficient control by imagining trajectories in latent space~\cite{hafner2019learning,hafner2020dream,ha2018world}.
Dreamer-style agents couple a recurrent state-space model (RSSM) with an actor--critic trained on imagined rollouts~\cite{hafner2020dream,hafner2023mastering}.
R2-Dreamer and related redundancy-reduced variants further streamline representation learning while retaining latent imagination~\cite{hafner2023mastering}.
Planners that search in latent space~\cite{hafner2019learning,okada2021dreaming} depend critically on transition accuracy; when the environment drifts, imagined rewards and values become biased even if the policy architecture is unchanged.

Most world-model agents, once trained, treat the dynamics module as \emph{fixed} at deployment.
MWM starts from the same pretrained backbone but asks: \emph{which parts of the model should remain frozen, which should adapt slowly with gradients, and which should adapt instantly without backprop?}
Our answer is to keep the bulk of $\theta$ fixed, add a separate fast state $P_t$ on prior logits, and optionally fine-tune a small slow subset of $\theta_{\mathrm{wm}}$ under consolidation.

%----------------------------------------------------------------------
\subsection{Test-time adaptation and online fine-tuning}
\label{sec:related-tta}

Classical meta-reinforcement learning learns initializations or update rules that can adapt quickly to new tasks from few gradient steps~\cite{finn2017model,duan2016rl}.
At deployment, adaptation still relies on backpropagation through policy or model parameters, often requiring careful inner-loop learning rates and task distributions matched at meta-train time.

\emph{Test-time training} and \emph{test-time adaptation} in supervised settings adjust batch-normalization statistics or encoder weights from unlabeled test batches~\cite{sun2020test,wang2021tent}.
Analogues in RL collect post-shift trajectories and continue policy or model training online~\cite{nagabandi2018learning,zintgraf2021fast,chua2018deep}.
RL$^2$ encodes adaptation in recurrent hidden state via sequence training~\cite{duan2016rl}, blurring the line between memory and gradient-based change but requiring meta-training across episodic variations.

These approaches share a \textbf{slow, global} adaptation mechanism: weight updates propagate through the network, are scheduled in batches, and risk catastrophic forgetting of pretraining~\cite{kirkpatrick2017overcoming}.
MWM is complementary: we add a \textbf{fast, local} store ($P_t$) that updates every step from prediction error without autograd through $\theta$, then use standard slow actor fine-tuning on replay for control.
Our experiments show that the fast pathway alone is insufficient for reliable return gains but materially improves slow actor learning when kept online (Section~\ref{sec:exp-bc})---a division of labor rather than a replacement for gradient TTA.

%----------------------------------------------------------------------
\subsection{Fast memory, context, and dual-process learning}
\label{sec:related-dual}

Several architectures separate transient context from persistent parameters.
Hypernetworks and fast weights modulate effective connectivity as a function of input or episode history~\cite{ha2017hypernetworks,ba2016using,mallya2018piggyback}.
Key--value episodic memory retrieves past transitions for one-shot or few-shot control~\cite{blundell2016model,pritzel2017neural}.
In sequence models, in-context learning can implement implicit adaptation without explicit gradient steps at test time~\cite{brown2020language}.

Cognitive science proposes \emph{complementary learning systems}: a fast hippocampal store for rapid encoding and a slow neocortical store for gradual integration~\cite{mcclelland1995complementary,kumaran2016learning}.
Recent machine-learning work revisits this split for continual and few-shot learning~\cite{ritter2018been,hung2019optimizing}.

MWM instantiates a CLS-like split \emph{inside} a world model:
$P_t$ is the fast, capacity-limited trace written by local error;
actor and $\theta_{\mathrm{wm}}$ are slow stores updated from replay;
coupled flush (Section~\ref{sec:consolidation}) is an engineered analogue of transferring hippocampal traces into cortical weights.
Unlike generic fast-weight layers, $P_t$ is tied explicitly to \textbf{RSSM prior prediction error} and norm-bounded Hebbian writes, making the fast store interpretable as a dynamics compensator rather than arbitrary context.

%----------------------------------------------------------------------
\subsection{Hebbian, local, and non-gradient plasticity}
\label{sec:related-hebbian}

Hebbian and covariance-based rules update synaptic efficacy from pre- and post-synaptic activity without global loss minimization~\cite{hebb1949organization,bienenstock1982theory}.
Modern work reintroduces local plasticity for continual learning, associative recall, and fast modulation in recurrent networks~\cite{ba2016using,mallya2018piggyback}.

Outer-product updates of the form $\Delta W \propto e \otimes x$ appear in linearized adaptation, fast-weight RNNs, and low-rank fine-tuning~\cite{hu2022lora}.
MWM's write rule (Eq.~\eqref{eq:plastic-write}) is in this family: signed error $e_t$ on prior logits and input $x_t$ from latent state and action jointly determine a rank-one increment to $P_t$, with decay and Frobenius projection ensuring stability.

Distinctions from low-rank adaptation (LoRA)~\cite{hu2022lora} and adapters~\cite{houlsby2019parameter}:
those methods introduce \emph{trainable} parameters at pretraining or fine-tuning time and update them by backprop;
$P_t$ has \textbf{no pretrained weights}, is \textbf{zeroed at episode boundaries} (optional), and is updated \textbf{only at test time} by a fixed Hebbian rule.
MWM operates in latent logit space of a categorical RSSM with explicit surprise gating and consolidation flush, targeting model-based RL under gain shift rather than supervised TTA or parameter-efficient fine-tuning alone.

%----------------------------------------------------------------------
\subsection{Adapting world models after deployment}
\label{sec:related-wm-adapt}

Prior work adapts dynamics models online for MPC~\cite{nagabandi2018learning,chua2018deep}, fine-tunes latent models under domain shift~\cite{hansen2022temporal}, or retrains decoders for sim-to-real~\cite{openai2018learning}.
Dreamer-style agents can in principle continue WM training from replay; in practice, unconstrained joint updates of dynamics and policy are unstable under non-stationarity, and imagination quality may degrade before control improves.

MWM constrains \emph{where} and \emph{how fast} each component may change:
fast prior correction ($P_t$), slow actor-only gradients with frozen-WM imagination (Section~\ref{sec:exp-actor-sweep}), and optional slow $\theta_{\mathrm{wm}}$ with sync to frozen inference copies.
Our WM-imagination ablation shows that plasticized multi-step rollouts during actor updates can \emph{hurt} performance even when online $P_t$ helps---supporting a split between \textbf{one-step belief correction} and \textbf{multi-step planning dynamics}.

%----------------------------------------------------------------------
\subsection{Positioning summary}
\label{sec:related-summary}

Table~\ref{tab:related-positioning} contrasts MWM with nearby paradigms.

\begin{table}[t]
  \centering
  \caption{Positioning of MWM relative to adjacent approaches.}
  \label{tab:related-positioning}
  \small
  \begin{tabular}{@{}lccc@{}}
    \toprule
    Approach & Fast store & Slow update & Targets WM prior \\
    \midrule
    Meta-RL / MAML~\cite{finn2017model} & -- & grad, all weights & optional \\
    Online fine-tuning~\cite{nagabandi2018learning,zintgraf2021fast} & -- & grad, batched & sometimes \\
    RL$^2$ / episodic memory~\cite{duan2016rl,pritzel2017neural} & hidden / KV & grad at meta-train & indirect \\
    LoRA / adapters~\cite{hu2022lora,houlsby2019parameter} & -- & grad, pretrained & any layer \\
    \textbf{MWM (ours)} & \textbf{$P_t$ Hebbian} & \textbf{actor / $\theta_{\mathrm{wm}}$} & \textbf{prior logits} \\
    \bottomrule
  \end{tabular}
\end{table}

In short, MWM is not primarily a new world-model architecture or a meta-learning algorithm.
It is a \textbf{test-time protocol} for frozen pretrained agents: augment RSSM priors with a norm-bounded Hebbian residual, pair it with slow replay-based actor learning, and optionally consolidate into $\theta_{\mathrm{wm}}$ while flushing $P_t$.
To our knowledge, this specific tri-timescale split---fast prior plasticity, frozen-imagination actor adapt, and coupled consolidation flush---has not been systematically evaluated in Dreamer-class agents under controlled actuator gain shift.
